\begin{abstract}
Out-of-distribution (OoD) inputs remain a major failure mode in modern object detectors, often leading to overconfident hallucinated detections on semantically irrelevant or unknown regions. While recent approaches rely primarily on post-hoc scoring functions or curated OoD datasets, they largely overlook the role of low-level signal characteristics in driving such failures.

In this work, we study how frequency bias and object scale contribute to objectness miscalibration in real-time object detectors. Motivated by spectral bias in deep neural networks, we show that high-frequency textures and small object regions disproportionately trigger false positive detections under OoD settings. To address this, we propose a frequency-aware objectness calibration framework that suppresses detection confidence using local frequency energy, objectness priors, and bounding box scale.

Our method is a lightweight post-hoc calibration mechanism that requires neither additional training data nor model modifications. We evaluate it using quantitative and explainability-driven metrics, including hallucination rate reduction, confidence redistribution, occlusion sensitivity, insertion-deletion faithfulness, and Grad-CAM-based localization consistency.

Experimental results show that our approach achieves up to 27.03\% reduction in false positive detections and substantially improves hallucination rates on OoD datasets while preserving in-distribution performance. Explainability analyses further show improved alignment between model attention and semantically meaningful regions after calibration.

Overall, this work highlights the importance of frequency-domain cues in OoD detection and provides a practical, low-compute strategy for mitigating hallucinations in real-time object detection systems.
\end{abstract}
